%! Composition of Functions — Unit 6, Lesson 6.6 — Exit Ticket
\documentclass[10pt]{article}
\usepackage{algebra2-article}
\usepackage{algebra2-boxes}
\newcommand{\TallMath}[1]{$\displaystyle #1\rule[-1.4em]{0pt}{3.2em}$}

\begin{document}

\pageheader{Unit 6, Lesson 6.6}{Exit Ticket}
\namedateperiod

\vspace{0.08in}

No notes. Work \textbf{inside out} --- and show the inside step, not just the final number.

\begin{enumerate}[leftmargin=1.8em, itemsep=10pt]
  \item \textbf{Both orders.} Let $f(x)=3x-2$ and $g(x)=x^{2}+1$.
        \vspace{0.62in}

        (a) $(f\circ g)(2)=\blank{1.8cm}$ \hspace{0.35in}
        (b) $(g\circ f)(2)=\blank{1.8cm}$

        \vspace{4pt}
        (c) Your two answers are different. Which function did you evaluate \emph{first} in part (a)?
        \blank{4.2cm}

  \item \textbf{Algebra and domain.} Let $f(x)=\sqrt{x}$ and $g(x)=x+7$.
        \vspace{0.55in}

        (a) $(f\circ g)(x)=\blank{3.0cm}$ \hspace{0.3in}
        (b) Domain: \blank{3.4cm}

        \vspace{5pt}
        (c) In one sentence, say why the domain is \emph{not} all real numbers --- your sentence
        should name which \emph{stage} of the composition does the refusing. \writeline

  \item \textbf{SOL-style multiple choice.} If $f(x)=x^{2}$ and $g(x)=x-5$, which expression is
        equal to $(f\circ g)(x)$?
        \begin{enumerate}[label=\Alph*., leftmargin=1.9em, itemsep=1pt]
          \item $x^{2}-5$
          \item $(x-5)^{2}$
          \item $x^{2}-25$
          \item $x^{3}-5x^{2}$
        \end{enumerate}
\end{enumerate}

\end{document}
